\documentclass[11pt,letterpaper]{article}
\usepackage[margin=1in]{geometry}
\usepackage[T1]{fontenc}
\usepackage{lmodern}
\usepackage{amsmath,amssymb,amsthm}
\usepackage{microtype}
\usepackage{booktabs,longtable,array}
\usepackage{enumitem}
\usepackage{xcolor}
\usepackage{xurl}
\usepackage{hyperref}
\usepackage{fancyhdr}
\definecolor{inkblue}{RGB}{27,55,78}
\hypersetup{colorlinks=true,linkcolor=inkblue,citecolor=inkblue,urlcolor=inkblue,
 pdftitle={Smaller explicit Diophantine forms of ZFC consistency},
 pdfauthor={AI-assisted research for the ProveIt repository}}
\newcommand{\Z}{\mathbb Z}
\newcommand{\N}{\mathbb N}
\DeclareMathOperator{\Pos}{Pos}
\newtheorem{theorem}{Theorem}[section]
\newtheorem{lemma}[theorem]{Lemma}
\newtheorem{proposition}[theorem]{Proposition}
\theoremstyle{remark}
\newtheorem{remark}[theorem]{Remark}
\setcounter{tocdepth}{2}
\setlength{\headheight}{14pt}
\setlength{\emergencystretch}{2em}
\pagestyle{fancy}
\fancyhf{}
\fancyhead[L]{\small\textsc{Diophantine forms of ZFC consistency}}
\fancyhead[R]{\small ProveIt synthesis}
\fancyfoot[C]{\thepage}
\renewcommand{\headrulewidth}{0.3pt}
\setlist{itemsep=3pt,topsep=5pt}
\allowdisplaybreaks[1]

\begin{document}
\hypersetup{pageanchor=false}
\begin{titlepage}
\vspace*{0.55in}
{\color{inkblue}\Large\textsc{ProveIt research report}\par}
\vspace{0.35in}
{\huge\bfseries Smaller explicit Diophantine\par
 forms of ZFC consistency\par}
\vspace{0.3in}
{\Large Twenty-four integers, tighter twelve-variable norms,\par
 and smaller materialized coefficients\par}
\vspace{0.35in}
{\large AI-assisted mathematical synthesis\par}
{\large Prepared for Vladimir Reshetnikov\par}
\vspace{0.15in}
{22 October 2026\par}
\vspace{0.4in}
\noindent\rule{\textwidth}{0.6pt}
\begin{abstract}
\noindent
Combining rational Pell saturation, separating-root norms, signed-domain
arguments and a concrete padding controller gives an explicit polynomial
in 24 unrestricted integers of degree 1936 with a 167-operation circuit.
ZFC proves that its lack of integer solutions is equivalent to
$\operatorname{Con}(\mathrm{ZFC})$.  Its three-equation version has degree
1920 and uses 162 operations.  Three twelve-positive-variable alternatives
have degree/operation pairs $(1069,123)$, $(1073,121)$ and $(1081,120)$.
An exact refactoring reduces the existing degree-1353 twelve-variable
equation to 113 operations.  Retaining the 943-tile language while
replacing five-power rounding reduces the fixed base to
$B=2^{6364399}$; all four resulting integers are materialized.
The report supplies explicit polynomial definitions, ordinary
mathematical proofs, arithmetic circuits and reproducible finite audits.
It does not claim a global minimum or a new kernel-checked end-to-end
formalization.
\end{abstract}
\vfill
\noindent\small
Source snapshot: \texttt{ae3007ca4467011e0e051c84acb4a1dd13971b0e}.\\
The archive preserves source provenance and separates new proofs from
inherited compiler and proof-calibration results.
\end{titlepage}

\hypersetup{pageanchor=true}
\pagenumbering{roman}
\tableofcontents
\clearpage
\pagenumbering{arabic}
% BEGIN sections/front.tex
\section{Results and the precise logical claim}
\label{sec:results}

There is no single useful meaning of ``simplest'' for these equations.
Number of unknowns, quantified domain, total polynomial degree, arithmetic
circuit length, and coefficient height are different resources.  This report
gives several improvements over the supplied collection, while preserving
its calibrated blank-machine construction.  The most substantial new
integer-domain result is a polynomial in \textbf{24 unrestricted integers},
of exact total degree \textbf{1936}, with a \textbf{167-operation} evaluator.
Its three-equation version has maximum degree 1920 and uses 162 operations.

For strictly positive integer unknowns, the report gives three new
twelve-variable equations with degree/operation pairs
\[
 (1069,123),\qquad (1073,121),\qquad (1081,120).
\]
An exact refactoring of the existing degree-1353 twelve-variable polynomial
also reduces its evaluator from 116 to 113 operations.  A separate
coefficient improvement reduces the radix exponent from 9,765,625 to
6,364,399, and supplies all four resulting fixed integers as actual binary
payloads.  These are complementary improvements, not a claim of a globally
minimal Diophantine representation.

\subsection{Comparison with the supplied snapshot}

All comparisons refer to the repository snapshot
\texttt{ae3007ca4467011e0e051c84acb4a1dd13971b0e}, rather than a moving
\texttt{main}.  The collection's own reconciliation identifies packages
16--23 as the current wave and records the historical alternatives
\cite{collection}.  Table~\ref{tab:comparison} compares like domains and
the particular resource being improved.

\begin{table}[htbp]
\centering\small
\begin{tabular}{p{2.31in}rrr}
\toprule
Single equation & Unknowns & Degree & Operations\\
\midrule
Previous integer, least count and degree (19) & $26\Z$ &2465&192\\
Previous integer, fewest ops at that count (21) &$26\Z$&2657&168\\
\textbf{New signed separating norm} &$\mathbf{24\Z}$&\textbf{1936}&\textbf{167}\\
\midrule
Previous positive, least degree at 12 (20)&12&1145&123\\
New tightest separating weight&12&1069&123\\
New intermediate weight&12&1073&121\\
New simpler weight&12&1081&120\\
\midrule
Previous positive, fewest ops at 12 (21)&12&1353&116\\
Same polynomial, refactored here&12&1353&113\\
\midrule
Previous monic merger (22)&13&319&91\\
Same polynomial, refactored here&13&319&89\\
\bottomrule
\end{tabular}
\caption{Positive integers unless marked $\Z$.  The integer comparison is
with the former minimum-variable family, not with every larger-variable
equation in the collection.  Source sizes: \cite{repo19,repo20,repo21,repo22}.}
\label{tab:comparison}
\end{table}

The new 24-integer equation improves both recorded minima of the former
26-integer family: its degree is below 2465 and its operation count is
below 168.  Larger-variable integer equations still have much lower
degrees; for example package 23 supplies 27 integers at degree 172
\cite{repo23}.  A lower degree is not being claimed across that tradeoff.

Two existing constructions also retain their particular advantages.
Package 18 gives nine natural-number unknowns at degree
$161844\cdot5^{33}+484$, with 20,531 operations; translating every unknown
by one gives nine positive unknowns at the same degree
\cite{repo18}.  The 15-positive-variable polynomial $P_{67}$ remains a
67-operation equation, of degree 155.  The retained-variable system from
package 22 has 16 positive unknowns, seven equations of maximum degree 24,
and 59 operations \cite{baseline67,retained59}.  Both of these latter
circuits are included with the new complete coefficient tuple.  Their
algebra and operation counts are unchanged.  Thus this report does not
claim fewer than nine positive unknowns, fewer than 67 operations for a
single equation, or fewer than 59 operations for a system.

\subsection{The theorem is an equivalence inside ZFC}

Fix the usual arithmetization of finite proofs in the chosen presentation
of ZFC, and let
\[
 \operatorname{Con}(\mathrm{ZFC})\equiv
 \neg\exists p\,\operatorname{Prf}_{\mathrm{ZFC}}
                         (p,\ulcorner\bot\urcorner).
\]
Here $\operatorname{Prf}_{\mathrm{ZFC}}$ is the finite syntactic proof
predicate.  The repository's proof predicate and any ordinary effective
presentation are related by the standard proof translations formalizable
in ZFC; no semantic soundness assumption is used in the equivalence.

Let $\boldsymbol\kappa=(\beta,K_0,M_C,M_F)$ be the four particular
integers defined in Section~\ref{sec:coefficients}.  For the 24-variable
polynomial $P_{24}$ explicitly defined in Section~\ref{sec:integer},
the claim is
\begin{equation}\label{eq:main-logical}
 \boxed{\quad\mathrm{ZFC}\ \vdash\
 \bigl(\neg\exists\boldsymbol x\in\Z^{24}\;
       P_{24}(\boldsymbol x;\boldsymbol\kappa)=0\bigr)
 \ \longleftrightarrow\ \operatorname{Con}(\mathrm{ZFC}).\quad}
\end{equation}
The analogous statements hold over $\Z_{>0}^{12}$ for each displayed
twelve-variable polynomial and over the corresponding domain for every
listed system.  Integer quantifiers here are an arithmetical abbreviation
inside ZFC.  The coefficients are fixed numerals; none is an extra
unknown, an unspecified oracle, or a value defined by waiting for the
machine to halt.

The construction proves the stronger positive equivalence
\begin{equation}\label{eq:existential-chain}
 \exists\boldsymbol x\;P(\boldsymbol x)=0
 \quad\longleftrightarrow\quad
 \operatorname{BlankHalts}(K215)
 \quad\longleftrightarrow\quad
 \exists p\;\operatorname{Prf}_{\mathrm{ZFC}}(p,\ulcorner\bot\urcorner).
\end{equation}
Negating this equivalence gives \eqref{eq:main-logical}.  In particular,
one does not assume $\operatorname{Con}(\mathrm{ZFC})$ while proving the
reduction, and this report does not prove the absence of solutions.
Under the usual hypotheses of the second incompleteness theorem, if ZFC
is consistent it cannot itself prove that absence.  Bare consistency
alone is not used to claim that ZFC cannot prove the opposite existential
sentence; that additional assertion would require a stronger soundness
hypothesis.

\subsection{What is inherited, and what is proved here}

The last equivalence in \eqref{eq:existential-chain} is the calibrated
K215 endpoint.  The supplied Lean source names it
\texttt{blank\_halts\_iff\_zfc\_refutation}, followed by
\texttt{nonhalting\_iff\_zfc\_consistency} \cite{K215}.  A byte-identical
copy and the machine table are included.  This delivery does not rebuild
that Lean development or represent the new number-theoretic reductions
as kernel-checked theorems.

The inherited native-window compiler represents accepting blank-machine
tableaux by Pell equations and fixed masks
\cite{baseline67,normalized,wang}.  Its soundness applies to arbitrary
solutions; its completeness constructs canonical witnesses from every
finite accepting computation.  The signed-domain and odd-scale results
specify which coordinates need positivity and how square and odd
substitutions preserve completeness \cite{normalized,oddscale}.
The source proof files needed for these interfaces accompany the report.

The new proofs supply three separate bridges.  First, rational Pell
saturation recovers an integer auxiliary root after root separation,
allowing a smaller multiplier than the previous separating construction.
Tighter weights then cancel the top homogeneous terms in four of the
sixteen linear conjugate factors.  Second, a signed separation argument,
including the $w=0$ and $T=0$ boundary cases, allows only six coordinates
to retain positivity guards; the resulting lift uses 24 integers.
Third, a concrete Boolean padding controller replaces the old
five-power rounding of the coefficient layout.  All three are ordinary
mathematical arguments using finite arithmetic, elementary algebraic
number theory, Pell identities and finite computation, hence formalizable
in ZFC.

\subsection{Conventions that keep the size claims meaningful}

The degree is the total degree of the actual polynomial after all domain
substitutions and after fixing the four coefficients.  The maximum
residual degree is reported for a system.  A circuit charges every binary
addition, subtraction and multiplication, including multiplication by a
fixed numeral; using a previously computed value and comparing two
values in a system are free.  Squaring is one multiplication.  The
coefficient numerals are input constants, so their bit lengths are
reported separately from circuit length.

Named intermediate expressions are polynomial abbreviations, not
additional existential variables.  Products over formal radicals below
are short descriptions of integer polynomials.  The delivered circuits
evaluate them using only $+,-,\times$, and the report also gives explicit
division-free formulas.  No root extraction or resultant is an uncharged
primitive.  Positive-domain results are always labeled as such; a
twelve-positive-variable equation is not asserted to have the same
soundness over twelve unrestricted integers.


% END sections/front.tex

% BEGIN sections/positive.tex
\section{Twelve positive unknowns: a tighter separating norm}
\label{sec:positive}

The construction combines the separating-root encoding of package~23 with
rational Pell saturation from packages~19--22. Saturation removes the extra
factor $H$ in package~23's strong multiplier. A sharper separating weight
then deliberately cancels leading terms in four conjugate factors. The
result improves package~20's degree-$1145$, $123$-operation equation
\cite{repo19,repo20,repo21,repo23}. All statements in this section use the
entire fixed coefficient tuple specified elsewhere in this report.

\subsection{An explicit polynomial}

The twelve supplied unknowns are
\[
 J,F,\alpha,z,h,i,T,s,w,\eta,\zeta,t\in\Z_{>0}.
\]
All other symbols are polynomial abbreviations. The fixed ports are $\beta,K_0,M_C,M_F$, with $\beta=B-1$ and $M_F$ the shifted high
mask. Define
\begin{align*}
 q&=\beta J+1,& C&=q-F-\alpha,& X&=wq,& Y&=sq^3,& E&=XY,\\
 k&=\eta+\zeta,&c&=kY+\eta,&a&=E+Y,&H&=4a+3,&\Delta&=a^2+H,\\
 \ell&=EkY,& A&=1+\ell(\ell+k),&G&=1+\Delta c^2,&b_0&=X+ac.
\end{align*}
The transport and index expressions are
\begin{align*}
 R&=((q-1)(q-F)+\alpha)(q^2-1)+(M_C+qM_F)J,\\
 N_k&=k-hE-R,\\
 N_t&=(K_0+w)C+q-F-z(q-1).
\end{align*}
The normalized strong and auxiliary expressions are
\begin{align*}
 I&=ic^2,& S&=\Delta I,&Q&=S^2,&D&=1+IS,&d&=Q-1,\\
 v&=cT,& b&=c+RD,&e&=Q(Dv^2+b^2)-1,&B_1&=2Qvb.
\end{align*}
In particular $\Delta=(a+2)^2-1$ and $D=1+\Delta I^2$.
Choose one of the following three weights:
\begin{align}
 M_{1081}&=1+8I(T+2REic),\label{eq:m1081}\\
 M_{1073}&=I\bigl(2T+R[E+3(X+Y)]ic\bigr),\label{eq:m1073}\\
 M_{1069}&=I\bigl(2T+R[E+2(X+Y)+sq^2]ic\bigr).
 \label{eq:m1069}
\end{align}
Set
\[
 m=HM,\qquad N=1+2m,\qquad W=Ht+Nb_0.
\]
For formal roots $\tau^2=A$, $g^2=G$, put
\begin{align}
 \mathcal F(\tau,g)
 &=\bigl[d(W-m\tau-Ng)^2-H^2e\bigr]^2-H^4B_1^2D,
 \label{eq:formalF}\\
 \Omega&=\prod_{\tau=\pm\sqrt A,\,g=\pm\sqrt G}
                   \mathcal F(\tau,g).\label{eq:omega}
\end{align}
The product is an integer polynomial, without additional unknowns.
The requested single equation is
\begin{equation}
 \boxed{\quad P=N_kN_t(1+\Omega)-1=0.\quad}
 \label{eq:positiveP}
\end{equation}
The corresponding three-equation system is
\begin{equation}
 N_k=1,\qquad N_t=1,\qquad \Omega=0.
 \label{eq:positiveSystem}
\end{equation}

For clarity, no algebraic-number or resultant operation is needed to
compute $\Omega$. The following scalar circuit uses only integer arithmetic:
\begin{align*}
 r_0&=W^2+N^2G,&\Theta&=4W^2N^2G,&a_0&=m^2A,\\
 C_0&=d(r_0+a_0)-H^2e,&k_f&=4d^2a_0,&K_c&=k_fr_0,\\
 Z_0&=H^4B_1^2D,&x_0&=C_0^2+\Theta d^2-K_c-Z_0,
   &x_1&=2C_0d-k_f,\\
 n_0&=x_0^2+\Theta x_1^2-4Z_0K_c,
   &n_1&=2x_0x_1-4Z_0k_f,\\
 \Omega&=n_0^2-\Theta n_1^2.
\end{align*}
To verify this identity, write $\theta=2WN g$, so $\theta^2=\Theta$.
After eliminating $\tau$, the conjugate factor is
$(C^2-K-Z_0)^2-4KZ_0$, where
$C=C_0+d\theta$ and $K=K_c+k_f\theta$.
Its even and odd parts are $n_0,n_1$, respectively.
The system compares $n_0^2=\Theta n_1^2$ directly, saving the final
subtraction. Shared intermediate values and the factorization of $M$ are
charged exactly as in the attached JSON circuits.

\begin{theorem}\label{thm:positive12}
For each weight in \eqref{eq:m1081}--\eqref{eq:m1069}, the positive
integer solvability of \eqref{eq:positiveP} or
\eqref{eq:positiveSystem} is equivalent to blank halting of the calibrated
K215 machine. Consequently ZFC proves that absence of a positive integer
solution is equivalent to $\operatorname{Con}(\mathrm{ZFC})$.
The exact sizes are as follows; every binary addition, subtraction, or
multiplication, including one by a fixed numeral, is charged.
\end{theorem}

\begin{center}
\begin{tabular}{lrrrr}
\hline
Weight & System degree & System operations & Equation degree & Operations\\
\hline
$M_{1081}$ &1072&115&1081&120\\
$M_{1073}$ &1064&116&1073&121\\
$M_{1069}$ &1060&118&1069&123\\
\hline
\end{tabular}
\end{center}

The sum-of-squares equation
$(N_k-1)^2+(N_t-1)^2+\Omega^2=0$ has degrees
$2144,2128,2120$ and operation counts $123,124,126$, respectively.
The theorem concerns ordinary mathematics formalizable in ZFC. It does
not assert a new kernel-checked end-to-end reduction or prove consistency.

\subsection{The integrality lemma that removes the extra multiplier}

\begin{lemma}[Rational Pell saturation]\label{lem:saturation}
Let $s_0\ge2$ and $V$ be integers. If $y\in\mathbb Q$ and
\[
 s_0^2V^2-(s_0^2-1)y^2=1,
\]
then $y\in\Z$.
\end{lemma}
\begin{proof}
Write $y=u/v_0$ in lowest terms with $v_0>0$. Its equation implies
$v_0^2\mid s_0^2-1$. The integer
$D_0=(s_0^2-1)/v_0^2$ is positive and nonsquare, since otherwise
$s_0^2-1$ would be a square. The pairs
$(s_0,v_0)$ and $(s_0|V|,|u|)$ solve $x^2-D_0z^2=1$ with positive
ordinates; $u=0$ or $V=0$ is impossible. Write them as
$(\chi_j,\psi_j)$ and $(\chi_l,\psi_l)$ in the powers of the fundamental
positive integer-coordinate Pell unit, where $j,l\ge1$.

We have $\chi_j\mid\chi_l$. The addition identities imply
\[
 \chi_{r+2j}\equiv-\chi_r\pmod{\chi_j},\qquad
 \chi_{2j-r}\equiv-\chi_r\pmod{\chi_j}.
\]
Reduce $l$ modulo $2j$, then reflect to $0\le r\le j$.
Strict growth gives $0<\chi_r<\chi_j$ for $r<j$, so the divisible
case requires $r=j$. Thus $l$ is an odd multiple of $j$.
The Pell multiplication formula gives $\psi_j\mid\psi_l$, hence
$v_0\mid u$. Their coprimality forces $v_0=1$.
\end{proof}

This is the special feature of the auxiliary norm. Rationality of a root
of a general nonmonic quadratic would not imply integrality. Packages
19--22 establish equivalent saturation lemmas
\cite{repo19,repo20,repo21}; applying the lemma before recovering the main
congruence removes package~23's extra $H$ without changing any canonical
compiler witness.

\subsection{Separation bounds before recovering roots}

First, $\Theta$ is divisible by $4$, so $\Omega$ is a square modulo $4$.
At a zero of \eqref{eq:positiveP}, its three integer factors are units.
The last cannot be $-1$: that would require $\Omega=-2$, which is not a
square modulo $4$. Therefore
\[
 \Omega=0,\qquad N_k=N_t=\epsilon\in\{1,-1\}.
\]
The system has $\epsilon=1$. The inherited weak transport bootstrap is
valid for both signs and gives, before any eliminated root is recovered,
\begin{equation}
 \begin{gathered}
 C\ge0,\quad q\ge16,\quad X\ge16,\quad Y\ge q^3\ge4096,\\
 0<R<q^4-2,\qquad E\ge q^4>R+2.
 \end{gathered}\label{eq:weakbounds}
\end{equation}
These are coefficient-dependent statements from the authentic blank
interface, not assertions for arbitrary substituted ports
\cite{normalized,repo20,repo21}.

A vanished factor in \eqref{eq:omega} supplies coherent roots satisfying
\begin{align}
 \tau^2&=A,&g^2&=G,&f^2&=D,\nonumber\\
 d y^2&=e-B_1f=Q(vf-b)^2-1,\label{eq:rootrelations}\\
 t&=y+M\tau+N\gamma,&\gamma&=(g-b_0)/H.
 \label{eq:linearroot}
\end{align}
The roots here are initially algebraic. In particular, no integrality
or sign of $\gamma$ has been assumed.

For the first weight, all complex auxiliary branches satisfy
\[
 |y|<2\bigl(v(a+2)I+b\bigr).
\]
Indeed $\sqrt D<(a+2)I$ and $Q\ge2$ give this bound directly from
\eqref{eq:rootrelations}. Since
$a+2<2E$, $c<2kY$, $\sqrt A>\ell>Ec/2$, and $D<4E^2I^2$,
\[
 \frac{|y|}{\sqrt A}
 <8TI+\frac4E+16REi^2c^3<M_{1081}.
\]

Here is the sharper estimate used for the other weights.
Since $\Delta>R+2$ and $I^2\ge c^4$,
\[
 Q-b=\Delta I^2(\Delta-R)-c-R>0.
\]
Thus $b\le d$; also $\sqrt D<S<d$.
Put $Z=v\sqrt D+b>1$. For every auxiliary conjugate,
\[
 |y|^2\le\frac{QZ^2+1}{d},\qquad
 \left[Z\left(1+\frac1d\right)\right]^2-
       \frac{QZ^2+1}{d}=\frac{QZ^2-d}{d^2}>0.
\]
Consequently
\begin{equation}
 |y|<Z+v+1<cT\bigl((a+2)I+1\bigr)+c+RD+1.
 \label{eq:sharpy}
\end{equation}
Let $\rho=(1+1/Y)/E$. We have $c/\ell<\rho$, and
$(a+3)\rho<2$, because
$(a+3)/E<9/8$ and $1+1/Y<17/16$.
Furthermore
\begin{align*}
 (a+2)^2/E&=E+2Y+4+Y/X+4/X+4/E<E+3Y+5,\\
 (a+2)^2\rho&<E+3Y+X+9<E+3(X+Y)-2.
\end{align*}
For $C_1=E+3(X+Y)$, dividing \eqref{eq:sharpy} by $\sqrt A$
therefore gives
\[
 \frac{|y|}{\sqrt A}
 <2TI+R(C_1-2)i^2c^3+2
 \le 2TI+RC_1i^2c^3=M_{1073}.
\]
The last inequality uses $Ri^2c^3\ge1$.

For $C_2=E+2(X+Y)+sq^2$, use the exact identity
\[
 (a+2)^2\rho=E+2Y+X+6+Y/X+5/X+4/Y+8/E+4/(EY).
\]
Since $Y/X=sq^2/w\le sq^2$ and
$5/X+4/Y+8/E+4/(EY)<1$, the right side is less than
$E+2Y+X+sq^2+7<C_2-2$. The identical argument proves
$|y|<M_{1069}\sqrt A$.

Finally $G>A$: $c>kY$ and $\Delta-E^2>2EY$ imply
\[
 \Delta c^2>k^2Y^2\Delta>
 k^2(E^2Y^2+EY)=A-1.
\]
For each of the three weights, the two uniform separation inequalities are
\begin{equation}
 |y|<M\sqrt A,\qquad
 \frac{N\sqrt G}{H}>2M\sqrt G>2M\sqrt A.
 \label{eq:separate}
\end{equation}

\subsection{Rationality, integrality, and the compiler interface}

Work in a finite normal splitting field over $\mathbb Q$ containing all
roots in \eqref{eq:rootrelations}. If an automorphism changed the sign of
$g$, it would change $N\gamma$ by magnitude $2N\sqrt G/H$.
The change in $M\tau$ has magnitude at most $2M\sqrt A$, and the change
in $y$ has magnitude strictly less than $2M\sqrt A$, using the uniform
bound for both conjugates. Their sum cannot cancel the former change by
\eqref{eq:separate}, contradicting the fixed rational value $t$ in
\eqref{eq:linearroot}. Hence $g$ is rational. With $g$ fixed, a change in
$\tau$ would similarly require $2M\sqrt A<2M\sqrt A$.
Thus $\tau$ is rational. Equation \eqref{eq:linearroot} makes $y$
rational, and $B_1>0$ makes $f=(e-dy^2)/B_1$ rational.

Since $g^2=G$, $\tau^2=A$, and $f^2=D$ are integers, $g,\tau,f$ are
integers. The integer $V=vf-b$ and $S=\Delta I\ge2$ now satisfy the
hypotheses of Lemma~\ref{lem:saturation}, so $y$ is an integer.
Both $H\gamma=g-b_0$ and $N\gamma=t-M\tau-y$ are integers, and
$N=1+2HM$ gives $\gcd(H,N)=1$. Bezout's identity proves
$\gamma\in\Z$.

The recovered equations are precisely the normalized first, main, strong,
and auxiliary Pell equations. The inherited signed-coordinate theorem
permits signed $\gamma,\tau,f,y,T$, and conversion to its ordinary strong
coordinate is $i_{\rm ord}=\Delta i$. That theorem restores the unit sign
$\epsilon=1$ and extracts a blank K215 halting computation
\cite{normalized,repo20,repo21}. In particular, the sign of the main root
is not assumed in the reconstruction.

Conversely, take a canonical positive normalized blank compiler witness.
For each chosen weight, set
\[
 t=y+M\tau+N\gamma>0.
\]
One factor of \eqref{eq:omega} vanishes and $N_k=N_t=1$.
No additional divisibility restriction on $i$ and no adjustment of the
auxiliary Pell period is required. This proves the existential equivalence
in Theorem~\ref{thm:positive12}; composing with the calibrated K215
refutation equivalence and negating existential solvability gives its
claimed ZFC theorem.

\subsection{Uniform exact degrees and the deliberate cancellation}

Scale the twelve supplied variables by a formal parameter $\lambda$ and
write a subscript $0$ for the leading homogeneous part. The shared degrees
are
\[
 \begin{array}{c|rrrrrrrrrrr}
 \text{expression}&E,H&Y&c&R&\Delta&I&Q,d&D&b&A,G&M,N\\
 \hline
 \text{degree}&6&4&5&4&12&11&46&34&38&22&27,33.
 \end{array}
\]
Put
\[
 \mathcal B=R_0E_0^2i^2c_0^4,\qquad
 \mathcal D=R_0E_0Y_0i^2c_0^4.
\]
These have degrees $38$ and $36$.
The two formal main-quotient roots solve
\[
 H^2\gamma^2+2Hb_0\gamma+(b_0^2-G)=0.
\]
The coefficient degrees are $12,17,16$, with
$(b_0^2-G)_{16}=-H_0c_0^2\ne0$. Thus the two roots have orders $5$ and
$-1$; the larger has leading term $-2\ell_0/H_0$.
Also each auxiliary root has expansion
\[
 y=\varepsilon(b-vf)+O(\lambda^{-8}),\qquad \varepsilon\in\{1,-1\},
\]
since $Q/(Q-1)=1+O(\lambda^{-46})$, while $vf$ has order $23$.
These expansions can be taken in a Laurent-series extension of the
rational function field; all leading radicands are squares there.

For $M_{1081}$, $M\sqrt A$ has leading term $16\mathcal B$ and
$N\gamma$ on its large branch has leading term $-64\mathcal B$.
The exact factorization is
\begin{equation}
 \Omega=d^8H^{16}
  \prod\bigl(t-y-M\tau-N\gamma\bigr),
 \label{eq:16factor}
\end{equation}
with sixteen coherent branches, counting multiplicity. Their leading
magnitudes are $15,17$ four times each, and $47,49,79,81$ twice each,
in units of $\mathcal B$. No leading term vanishes. Hence
\[
 \deg\Omega=8\cdot46+16\cdot6+16\cdot38=1072.
\]
The leading coefficient multiplier is
$15^4 17^4(47\cdot49\cdot79\cdot81)^2$, which is nonzero.

For $M_{1073}$, the positive formal first root satisfies
$\sqrt A-\ell=O(\lambda)$, and
\[
 M\sqrt A-b=Ri^2c^3(C_1\ell-\Delta c)+O(\lambda^{28}).
\]
Here the exact polynomial identity
\[
 C_1\ell-\Delta c
 =(EY+3EX-Y^2-4E-4Y-3)c-E(E+3X+3Y)\eta
\]
has unique leading form $E_0Y_0c_0$ of degree $15$.
Consequently $M\sqrt A-b$ has leading form $\mathcal D$ of degree $36$:
its terms of degrees $38$ and $37$ cancel identically.
Now $N\gamma$ on the large branch has leading term $-4\mathcal B$;
on the small branch it has order $32$.
In \eqref{eq:16factor}, four factors have degree $36$ and leading
magnitude $\mathcal D$, four have degree $38$ and magnitude
$2\mathcal B$, and the remaining eight have leading multiples
$2,4,4,6$ of $\mathcal B$, each twice. The signed product is positive:
\begin{align*}
 \Omega_{1064}
 &=589824\,d_0^8H_0^{16}\mathcal B^{12}\mathcal D^4,\\
 1064&=8\cdot46+16\cdot6+12\cdot38+4\cdot36.
\end{align*}
Writing $q_0=\beta J$, $k_0=\eta+\zeta$ gives the fully substituted form
\begin{equation}
 \Omega_{1064}=9\cdot2^{48}i^{48}k_0^{96}w^{76}s^{176}
                     q_0^{652}(q_0-F)^{16}.\label{eq:top1073}
\end{equation}
It is nonzero for every authentic slice $\beta>0$.
The product's leading form multiplies this by
\[
 [-hw s q_0^4]\,[w(q_0-F-\alpha)-zq_0].
\]
For example, the coefficient of
$J^{673}hz i^{48}\eta^{96}w^{77}s^{177}$ is
$9\cdot2^{48}\beta^{673}>0$. This proves exact degree $1073$ uniformly,
not merely for diagnostic numerical specializations.

Finally, for $M_{1069}$ replace $C_1$ by $C_2$ above. Then
\[
 C_2\ell-\Delta c
 =(2EX+Esq^2-Y^2-4E-4Y-3)c-EC_2\eta.
\]
Its unique leading term is $E_0(sq_0^2)c_0$, of degree $14$.
Thus the four cancellation factors have degree $35$, rather than $36$.
All other factors retain their previous degrees and leading multipliers.
The norm has exact degree $1060$, and its leading form is
\[
 \Omega_{1060}=9\cdot2^{48}i^{48}k_0^{96}w^{76}s^{176}
                     q_0^{648}(q_0-F)^{16}.
\]
The corresponding product coefficient at
$J^{669}hz i^{48}\eta^{96}w^{77}s^{177}$ is
$9\cdot2^{48}\beta^{669}>0$, proving exact degree $1069$.
The degrees of $N_k,N_t$ are $7,2$ in every case. The sum-of-squares
claims follow because $\Omega^2$ is its unique highest-degree summand.

The independent finite checks compare the literal circuits against
conjugate-product evaluators, verify all gates and ports are live, and
compute full univariate specializations over two finite fields to detect
the cancellations. They corroborate the formulas and operation counts;
the preceding uniform arguments establish the mathematical equivalence
and exact degrees. They do not construct an accepting K215 computation.

% END sections/positive.tex

% BEGIN sections/integer.tex
\section{Twenty-four unrestricted integer unknowns}
\label{sec:integer}

The preceding construction can be adapted to a substantially smaller
integer domain. The current collection has equations in 26 unrestricted
integers; the construction below has \textbf{24}, degree \textbf{1936},
and a division-free evaluator with \textbf{167} binary operations. Its
three-equation version uses 162 operations and has maximum degree 1920.
The proof combines the saturated root norm of package~19 with the odd-scale
signed-domain theorem of package~23 \cite{repo19,repo23,oddscale}.
The new work is the separation argument when the half-binomial scale, the
main Pell parameter, and the separating coefficients can be negative.

\subsection{Domain maps and the equation}

For each of the six ports
\[
 J,F,\alpha,z,\eta,\zeta
\]
use a disjoint copy of the polynomial
\[
 \operatorname{Pos}(x,y,u)=1+x^2+y^2+u(u+1).
\]
Its values are positive, and it represents every positive integer:
apply the three-square theorem to $4n-3$, whose representation has one
odd square and two even squares. Supply six further unrestricted integers
$h,T,t,r_i,r_s,r_w$, and put
\begin{equation}\label{eq:integer-domain}
 i=2r_i+1,\qquad s=2r_s+1,\qquad w=r_w^2.
\end{equation}
Thus the number of supplied integers is $6\cdot3+6=24$.
The executable circuit names these last three inputs \texttt{ii},
\texttt{ss}, and \texttt{ww}, and names the merged coordinate $t$
\texttt{r}.

Retain the common abbreviations $q,C,X,Y,E,k,c,a,H,\Delta,R,N_k,N_t$
of the positive construction. In particular,
\begin{align*}
 X&=wq,&Y&=sq^3,&E&=XY,&k&=\eta+\zeta,\\
 c&=kY+\eta,&a&=E+Y,&H&=4a+3,&\Delta&=a^2+H.
\end{align*}
To make the norm definition explicit, set
\begin{align*}
 L&=kXY^2,& A&=1+L(L+k),&G&=1+\Delta c^2,\\
 P&=X+ac,&I&=ic^2,&D&=1+\Delta I^2,\\
 Q&=\Delta^2I^2,&d&=Q-1,&v&=cT,\\
 b&=c+RD,&e&=Q(Dv^2+b^2)-1,&B_1&=2Qvb.
\end{align*}
Here $L$ is an arithmetic abbreviation, unrelated to a compiler cell
length. All these symbols are computed expressions, not new unknowns.

Use the following signed separating weights:
\begin{align}
 M&=1+8(TI)^2+16Ra i^2c^3
     =1+8(ic)I(T^2c+2Ra),\label{eq:signed-weight}\\
 N&=1+8HM,\qquad Z=Ht+NP.\label{eq:signed-weight-N}
\end{align}
The second expression for $M$ is the factored form used by the circuit.
For formal roots $u^2=A$ and $g^2=G$, define
\begin{align}
 F(u,g)&=\bigl[d(Z-HMu-Ng)^2-H^2e\bigr]^2-H^4B_1^2D,
 \label{eq:signed-factor}\\
 \Omega&=\prod_{u=\pm\sqrt A}\prod_{g=\pm\sqrt G}F(u,g).
 \label{eq:signed-norm}
\end{align}
The product is an ordinary polynomial with integer coefficients; its
division-free circuit does not use square roots or a resultant operation.
The requested single equation is
\begin{equation}\label{eq:integer-final}
 \boxed{\;N_kN_t(1+\Omega)-1=0.\;}
\end{equation}
The corresponding system is $N_k=1$, $N_t=1$, $\Omega=0$.
All four fixed coefficient ports must come from one complete authentic
tuple. The new tuple in Section~\ref{sec:coefficients} has the required
interface; package~19's compact Wang943 tuple is a separate valid choice.

\begin{theorem}\label{thm:integer24}
For the coherent fixed coefficients specified in this report, equation
\eqref{eq:integer-final} has an integer zero if and only if the blank
run of K215 halts. The same statement holds for its three-equation system.
Consequently ZFC proves that the absence of integer zeros is equivalent
to $\operatorname{Con}(\mathrm{ZFC})$.
\end{theorem}

\subsection{Bootstrap with signed scales}

The optimized norm evaluator has the form
$\Omega=U^2-\Theta V^2$, where
$\Theta=4Z^2N^2G$. Thus $\Omega$ is a square modulo four.
At a zero of \eqref{eq:integer-final}, all three integer factors are
units. The alternative $1+\Omega=-1$ would require a square to equal
two modulo four. Therefore
\[
 \Omega=0,\qquad N_k=N_t=\epsilon,\qquad\epsilon\in\{1,-1\}.
\]
For the system these conclusions hold directly with $\epsilon=1$.

The weak transport estimate needs only the six guarded positive ports:
\[
 q-F-z(q-1)=1-F-(z-1)(q-1)\leq0.
\]
If $C\leq-1$, then $N_t\leq-(K_0+w)\leq-2$, a contradiction.
Hence $C\geq0$ and $q-F=C+\alpha>0$. The authentic mask bounds give
\begin{equation}\label{eq:integer-mask-bound}
 q\geq16,\qquad q(q^2-1)<R<q^4-2.
\end{equation}
These estimates are valid before reconstructing any root and do not use
$s>0$ \cite{normalized,repo19}.

Since $s$ is odd, $|Y|\geq q^3$. Also $X\geq0$, and $a,c$ have
the same sign as $Y$. Their absolute values satisfy
$|a|\geq q^3$ and $|c|>k(|Y|-1)>1$. Consequently $|H|>1$ and
\[
 \Delta=(a+2)^2-1\geq(q^3-2)^2-1>q^4>R.
\]
The multiplier $i$ is nonzero, so
\begin{equation}\label{eq:integer-pre-root}
 D>\Delta+|c|,\qquad 1<b<Q,\qquad Q>D>1.
\end{equation}
For the upper bound on $b$, use
$Q-b=(\Delta-R)D-\Delta-c>0$. In particular $b$ is positive even
when $c$ is negative. Finally, $Ra c^3>0$, so $M>1$, and
\begin{equation}\label{eq:integer-N-bound}
 \left|\frac NH\right|>7M,\qquad \gcd(N,H)=1.
\end{equation}

\subsection{Recovering roots without assuming their signs}

Because $A\geq1$ and $G>1$, a zero of \eqref{eq:signed-norm}
selects real roots $u,g$ with $F(u,g)=0$. Work in the totally real
Galois field $\mathbb Q(u,g)$, allowing degrees one, two, or four, and put
\[
 y=\frac{Z-HMu-Ng}{H}.
\]
When $B_1\ne0$, define $f=(e-dy^2)/B_1$. Then $f^2=D$ and every
actual conjugate satisfies
$d(y')^2=Q(vf'-b)^2-1$. When $B_1=0$, necessarily $v=0$, and the
same factor directly gives $d(y')^2=Qb^2-1$. In both cases,
\begin{equation}\label{eq:integer-y-bound}
 |y'|<2\bigl(|v|\sqrt D+b\bigr),\qquad
 Z=HMu+Ng+Hy.
\end{equation}
No division by $T$ or $B_1$ is made in the exceptional branch.

First exclude $w=0$. The inequalities $|a|\geq4096$ and $|c|>1$ give
\[
 \sqrt G>\frac{|a||c|}{2},\qquad
 \sqrt D<2|a||i|c^2,\qquad D<4a^2i^2c^4.
\]
Term by term in \eqref{eq:signed-weight}, these imply
\begin{align}
 M\sqrt G
 &>2|c|+4|a||iT||c|^3+8Ra^2i^2c^4\notag\\
 &>2\bigl(|v|\sqrt D+b\bigr).
 \label{eq:integer-boundary-separation}
\end{align}
Here $i^2T^2c^2\geq|iT|$, including $T=0$, and
$b\leq|c|+RD$. If $w=0$, then $A=1$ and $u$ is already integral.
An automorphism flipping $g$ would contradict
\eqref{eq:integer-N-bound}, \eqref{eq:integer-y-bound}, and
\eqref{eq:integer-boundary-separation}. Thus $g$ is rational and
integral, and $y$ is rational.

We use the rational auxiliary Pell saturation lemma: for integers
$S_0>1,V$, a rational solution of
\begin{equation}\label{eq:integer-saturation}
 S_0^2V^2-(S_0^2-1)y^2=1
\end{equation}
is integral. Its proof by ordinary Pell classification and divisibility
of first coordinates does not assume $S_0^2-1$ squarefree
\cite{repo19}. When $B_1\ne0$, the rational $f$ has integral square,
so is integral, and the lemma applies with
$S_0=|\Delta I|$, $V=vf-b$. When $B_1=0$, apply it with $V=-b$.
Thus $y$ is integral in either case.

Set $\gamma=(g-P)/H$. The linear equation gives
$t=Mu+N\gamma+y$. Hence $H\gamma$ and $N\gamma$ are integers;
coprimality in \eqref{eq:integer-N-bound} makes $\gamma$ integral.
At $X=0$ the main norm would now imply
\[
 1=(ac+H\gamma)^2-\Delta c^2
   =H(2ac\gamma+H\gamma^2-c^2),
\]
contradicting $|H|>1$. Therefore $w>0$ and $X\geq q\geq16$.

For $w>0$, we have $\sqrt A>L$, $|a+2|<2|E|$, and
$|c|<2k|Y|$. Dividing \eqref{eq:integer-y-bound} by $\sqrt A$
and using $L=|E|k|Y|$ yields
\begin{align}
 \frac{|y'|}{\sqrt A}
 &<8|TI|+\frac4{|E|}+16R|E|i^2|c|^3\notag\\
 &<1+8(TI)^2+16R|a|i^2|c|^3=M.
 \label{eq:integer-ordinary-separation}
\end{align}
The last step uses $|E|\geq q^4$, $|a|>|E|$, and the integer
inequality $(TI)^2\geq|TI|$.

For either sign of $Y$, the elementary estimates
\[
 \Delta>\tfrac34a^2,\quad |c|>\tfrac34k|Y|,
 \quad G>\tfrac{27}{64}L^2,\quad A<\tfrac32L^2
\]
give $\sqrt G>\sqrt A/2$. An automorphism flipping $g$ in
\eqref{eq:integer-y-bound} would therefore imply
$|N/H|\sqrt G<2M\sqrt A$, whereas
\eqref{eq:integer-N-bound} gives
$|N/H|\sqrt G>\tfrac72M\sqrt A$. Thus $g$ is rational. Once $g$
is fixed, a flip of $u$ contradicts the uniform bound
\eqref{eq:integer-ordinary-separation}. Hence $u$ is rational as well.
Both are integers, and the saturation and coprimality argument above
makes $y,\gamma$ integers. This uses only actual field automorphisms;
square radicands or coincident quadratic extensions cause no exception.

If $T=0$, the recovered integral $y$ would satisfy
\[
 (Q-1)(y^2-b^2)=b^2-1>0.
\]
Since $1<b<Q$, one has $|y|\geq b+1$ and $Q-1\geq b$. The left
side is at least $b(2b+1)>b^2-1$, a contradiction. Thus $T\ne0$,
and the recovered $f=(e-dy^2)/B_1$ is an integer.

\subsection{Soundness and completeness of the signed domain}

The roots just reconstructed satisfy exactly
\begin{align*}
 u^2&=1+L(L+k),\\
 (X+ac+H\gamma)^2-\Delta c^2&=1,\\
 f^2-\Delta i^2c^4&=1,\\
 Q[c(Tf-1)-Rf^2]^2-(Q-1)y^2&=1,
 \qquad N_k=N_t=\epsilon.
\end{align*}
There are six positive guarded ports, odd $s,i$, nonzero square $w$,
and unrestricted integer roots. These are precisely the hypotheses of
the signed odd-scale theorem \cite{oddscale}. Its negative-$s$ argument
recovers the main Pell rank $R$ and rules out both main-root signs by
their congruences modulo $H$. Hence $s>0$. The inherited signed compiler
then forces $\epsilon=1$ and recovers an accepting blank K215 tableau
\cite{normalized}. Our root recovery did not assume these conclusions.

Conversely, an accepting computation has canonical normalized witnesses
with odd $i$ and square $w$. The additional oddness of $s$ follows from
the native population identities. If $q=2^j$, the native masks give
\[
 \operatorname{popcount}(R)=3j+2,
 \qquad v_2(Y)=\operatorname{popcount}(R)-2=3j.
\]
Since $Y=sq^3$, the scale $s$ is odd. The required mask hypotheses are
$M_C\equiv2\pmod4$, $M_F^{\rm native}\equiv4\pmod8$, and
$\operatorname{popcount}(M_C)+
\operatorname{popcount}(M_F^{\rm native})=d_{\rm base}=\log_2B$;
the coefficient adapter preserves them. The arithmetic port is still
the \emph{shifted} $M_F$.

Choose $t=y+Mu+N\gamma$. Then one norm factor vanishes and
$N_k=N_t=1$. The six positive coordinates admit their three-integer
representations. This proves Theorem~\ref{thm:integer24} after composition
with the inherited K215/refutation calibration.

\subsection{Exact degree and operation count}

Assign weight two to $J,F,\alpha,z,\eta,\zeta,w$ and weight one to
$s,i,h,T,t$. The principal degrees are
\begin{align*}
 \deg(a,H,c,\Delta,L,A,G)&=(11,11,9,22,20,40,40),\\
 \deg(R,I,D,Q,d,b,e,B_1)&=(8,19,60,82,82,68,218,160),\\
 \deg(M,N,Z,N_k,N_t)&=(48,59,79,12,4).
\end{align*}
Write a subscript $0$ for the highest homogeneous part. Then
$A_0=G_0=L_0^2$, $P_0=L_0=a_0c_0$, and
\[
 M_0L_0=16b_0,\qquad N_0=8H_0M_0.
\]
In the four leading sign choices, $Z-HMu-Ng$ has leading part
$H_0M_0L_0$ times one of $-1,1,15,17$. Thus the leading coefficients
of the inner bracket in \eqref{eq:signed-factor}, after extracting
$H_0^2d_0b_0^2$, are
\[
 255,\quad255,\quad57599,\quad73983.
\]
Their squares have degree 480, while the competing term $H^4B_1^2D$
has degree 424. Therefore
\begin{equation}\label{eq:integer-top-form}
 \Omega_0=C_*H_0^{16}d_0^8b_0^{16},\qquad
 C_*=255^4(57599\cdot73983)^2>0.
\end{equation}
Its degree is $16\cdot11+8\cdot82+16\cdot68=1920$.
Multiplying $N_kN_t$ adds $12+4$, giving exact degree 1936.
The squared-residual alternative has degree 3840.

For an explicit uniform certificate, in the actual circuit variable
names the coefficient of
\[
 J_x^{1354}h\,\mathrm{ii}^{48}\eta_x^{192}
 \mathrm{ww}^{162}\mathrm{ss}^{177}z_x^2
\]
is $2^{257}C_*\beta^{677}$. The exponent sum is 1936 and the
coefficient is nonzero for every authentic $\beta>0$. Thus the degree
claim applies to the full fixed integers, independently of numerical
specializations.

The optimized positive core takes 120 product operations; the signed
weight takes two more. The domain maps initially cost 47 operations.
Two additions disappear by absorbing the constant terms of $F$ and $z$:
if $F=1+F_*$, $z=1+z_*$, and $r=\beta J$, then
\[
 q-F=r-F_*,\qquad N_t=(K_0+w)C-F_*-z_*r.
\]
The final counts are therefore 167 operations for the product, 162 for
the system, and 170 for the sum of squares. All supplied coordinates and
all emitted gates are live. Exact conjugate-product checks, domain-map
identity checks, and independent finite-field leading-coefficient checks
corroborate the circuit. The universal equivalence rests on the proof
above and the explicitly stated inherited dependencies, rather than on
those finite diagnostics.

% END sections/integer.tex

% BEGIN sections/circuits.tex
\section{Division-free circuits and arithmetic savings}
\label{sec:circuits}

All displayed auxiliary symbols in this section denote polynomial
abbreviations. They introduce neither existentially quantified coordinates
nor operations outside addition, subtraction, and multiplication.
The delivered JSON files specify their evaluation order and every operand.
The accounting convention charges each binary arithmetic operation, including
multiplication by a fixed integer; reusing an earlier value and comparing two
computed expressions are free. This is the convention used by the repository
comparison circuits \cite{repo21,repo22}.
All polynomial-identity comparisons hold the four coefficient ports equal.
Replacing the whole tuple is a separate semantic step proved in
Section~\ref{sec:coefficients}.

\subsection{A shared evaluator for the four-root construction}

Use the notation of the root-recovery construction, and write $W$ for its
supplied merged expression, $N$ for the weight of the main root, and $m=HM$
for the weight of the first Pell root. Thus the relevant eight-conjugate
product is
\begin{equation}
\mathcal R=
\prod_{\sigma,\epsilon,\delta\in\{-1,1\}}
\left[d\bigl(W+\sigma N\sqrt G+\epsilon m\sqrt A\bigr)^2
      -H^2e+\delta H^2B_{\rm aux}\sqrt D\right].
\label{eq:circuits-eight-norm}
\end{equation}
Here $B_{\rm aux}=2Qvb$. Define $P_a=QDv^2$ and $U_a=Qb^2$,
so that $B_{\rm aux}^2D=4P_aU_a$. Formula \eqref{eq:circuits-eight-norm} explains the elimination;
it is not the circuit definition and does not require a square-root operation.

The following completely integral evaluator computes the same polynomial:
\begin{align*}
 V_0&=W^2+N^2G,& J_\theta&=4W^2(N^2G),& a_0&=m^2A,\\
 C_0&=d(V_0+a_0)-H^2e,& K_*&=4d^2a_0,& U_0&=K_*V_0,\\
 Z_0&=4H^4P_aU_a,\\
 x&=C_0^2+J_\theta d^2-U_0-Z_0,
 &y&=2C_0d-K_*,\\
 n_0&=x^2+J_\theta y^2-4Z_0U_0,
 &n_1&=2xy-4Z_0K_*,\\
 \mathcal R&=n_0^2-J_\theta n_1^2.
\end{align*}
The three-equation system compares
\[
 N_k=1,\qquad N_t=1,\qquad n_0^2=J_\theta n_1^2.
\]
The corresponding single equation is
\[
 N_kN_t(1+\mathcal R)-1=0.
\]
Since $4\mid J_\theta$, the integer $\mathcal R$ is a square modulo $4$.
Consequently $1+\mathcal R$ cannot equal $-1$. This preserves the modular
guard needed in the product equation; the remaining outer signs are handled
by the root-recovery proof.

To verify the evaluator, work temporarily in a quadratic algebra and put
$\theta=2WN\sqrt G$, so that $\theta^2=J_\theta$ and
$(W+N\sqrt G)^2=V_0+\theta$. Taking the norm in the first-root direction
uses the exact identity
\begin{equation}
 \bigl((C-u)^2-Z\bigr)\bigl((C+u)^2-Z\bigr)
   =(C^2-u^2-Z)^2-4u^2Z.
\label{eq:circuits-symmetric}
\end{equation}
After the first two eliminations, the remaining quadratic-algebra element
is $n_0+n_1\theta$. Its norm is the last line of the evaluator. This is an
identity in a free integer-polynomial ring. In particular no division by
$W$, $N$, or $2WN$ is performed, and the identity remains valid when one of
these quantities is zero.

\subsection{Factoring the common odd coefficient}

The earlier paired evaluator explicitly constructs $V_1=2WN$, then
$dV_1$, $K_*V_1$, and their products and squares. Every odd component has
this same factor $V_1$. Retaining only its square
$G V_1^2=4W^2(N^2G)$ avoids three arithmetic operations in the complete
polynomial evaluator. Both $W^2$ and $N^2G$ are already needed for $V_0$.
The system saves one further subtraction by comparing its two final norm
terms directly.

Applied to the unchanged nested-norm polynomial of package~21, this gives
\[
 (12\hbox{ positive variables},\ 1\hbox{ equation},\ 1353,\ 113)
\]
in the order (unknowns, equations, degree, operations), improving its
116-operation evaluator. Its three-equation system uses 108 operations
rather than 112, with the same maximum degree 1344. The polynomial itself,
its zero set, and the proof assumptions are unchanged \cite{repo21}.

The same idea applies to the monic-resultant circuits of package~22.
There $r_1=2W$ occurs in every odd component, while $Br_1^2$ already occurs
in the even component. In its twelve-variable evaluator, write
\[
 v_I=r_1t,\qquad
 h_I=2r_1(v_Rt-\lambda u_R).
\]
One may compute $Bv_I^2=(Br_1^2)t^2$ and
$Bh_I^2=4(Br_1^2)(v_Rt-\lambda u_R)^2$ directly. Moreover
$Br_1^2=4BW^2$ needs two gates because $W^2$ is already available,
whereas forming $r_1$, squaring it, and multiplying by $B$ uses three.
The complete saving is three operations: 117 becomes 114 for the
polynomial, and 112 becomes 109 for its system. The thirteen-variable
intermediate has $h_I=r_1t$ and saves two operations: its degree-319
polynomial uses 89 operations, and its four-equation system of maximum
degree 288 uses 81. Again these are exact refactorings of the original
polynomials \cite{repo22}.

\subsection{Weights and the unrestricted-integer lift}

The new positive constructions use the same shared evaluator with different
separating weights. For the degree-1081 version, the identity
\[
 1+8TI+16REi^2c^3=1+8I(T+2REic),\qquad I=ic^2,
\]
saves one more operation. The tighter degree-1073 and degree-1069 weights
are evaluated in their already factored form
\[
 M=I(2T+R\mathcal C\,ic),
\]
with $\mathcal C=E+3(X+Y)$ and
$\mathcal C=E+2(X+Y)+sq^2$, respectively. Their improvement in degree
comes from the cancellation proved in the weight analysis, rather than
from the generic degree bound of the evaluation graph.

For the signed construction use
\[
 M=1+8(TI)^2+16Ra i^2c^3
   =1+8(ic)I(T^2c+2Ra),\qquad N=1+8HM.
\]
The second expression is a polynomial identity for arbitrary integers,
including negative $i$ and $c$. The separate signed-domain proof supplies
the positivity properties needed in the root argument.

There is a further saving when the six guarded coordinates are replaced by
$\Pos(u,v,t)=1+u^2+v^2+t(t+1)$. Write
\[
 F=F_0+1,\quad z=z_0+1,\quad q=\beta J+1,
\]
where $F_0$ and $z_0$ are their three-term nonnegative quadratic parts.
Then
\begin{align*}
 q-F&=\beta J-F_0,\\
 (K_0+w)C+(q-F)-z(q-1)
   &=(K_0+w)C-F_0-z_0\beta J.
\end{align*}
Thus the evaluator never forms the two separate $+1$ gates for $F$ and
$z$. This changes neither the substitution nor the quantified domain.
The final unrestricted-integer polynomial uses 167 operations, and its
three-equation system uses 162.

\subsection{Final ledgers and verification scope}

\begin{center}
\small
\begin{tabular}{lrrrr}
\hline
Form & Polynomial degree & Operations & System degree & Operations\\
\hline
12 positive, tightest weight &1069&123&1060&118\\
12 positive, tight weight    &1073&121&1064&116\\
12 positive, simpler weight &1081&120&1072&115\\
12 positive, refactored norm&1353&113&1344&108\\
13 positive, monic merger   &319&89&288&81\\
24 unrestricted integers    &1936&167&1920&162\\
\hline
\end{tabular}
\end{center}
All systems in this table have three equations except the thirteen-variable
system, which has four. The degrees are exact. For the degree-1069 and
1073 products, naive propagation through the graph gives 1081; their
homogeneous cancellation must be used to obtain the exact degree.
The graph metadata records this distinction rather than overwriting the
larger syntactic bound.

The package contains three complementary kinds of evidence. First,
\path{code/check_exact_identities.py} expands the small free-polynomial
identities over $\Z$ and checks every coefficient. This covers the shared
odd-component factors, identity~\eqref{eq:circuits-symmetric}, the weight
factorizations, and both integer-lift cancellations. Second, independent
exact evaluations multiply all eight conjugates in an eight-dimensional
algebra and compare that result with the division-free graph. Third,
\texttt{degree\_check.py} evaluates the entire graph on univariate slices
over two prime fields. Its convolution uses integer packing with a proved
carry bound; it does not use floating point or merely propagate the
largest input degree. These checks confirm the degree cancellations and
support the separate uniform leading-form proofs.

For the signed polynomial, the independently derived monomial
\[
 J_x^{1354}h\,i_{\rm int}^{48}\eta_x^{192}
 w_{\rm int}^{162}s_{\rm int}^{177}z_x^2
\]
has total degree 1936 and coefficient
\[
 2^{257}\,255^4(57599\cdot73983)^2\beta^{677}.
\]
The coefficient is nonzero for every admissible $\beta>0$; its modular
values agree with two independent full-polynomial slice calculations.

Arithmetic identity checks establish the stated circuit transformations.
The root-recovery lemmas, the signed-domain argument, the native compiler,
and the calibration of K215 remain mathematical proof dependencies. None
of these finite calculations is presented as a new Lean or Rocq theorem,
or as a test of whether the consistency polynomial has a solution.

% END sections/circuits.tex

% BEGIN sections/coefficients.tex
\section{Smaller materialized coefficients without prime-power rounding}
\label{sec:coefficients}

The new arithmetic forms can use the complete Wang943 coefficient tuple of
package 19.  A further improvement is possible by combining that package's
exact support calculation with the repository's Boolean padding lemma for
an arbitrary fixed odd factor \cite{wang,padding}.  The geometric radix
dimensions need not themselves be powers of five.  What completeness
actually requires is a suitable family of tableau sizes and enough freely
selectable dummy bits to tune the actual packed index.

\subsection{The literal finite recipe}

Retain all 943 directed-signal Wang tiles, their order, the occupied ten-bit
edge codes, the four copy tracks, and the native exponent positions of the
compact construction.  The first tile is Start, the second is the globally
impossible forbidden tile, and the remaining tiles are lexicographically
ordered.  The included 430-action machine table and tile list are checked
against one another by an independent enumeration in
\texttt{code/build\_coefficients.py}.  The machine table's SHA-256 is
\begin{center}\small\ttfamily
fdbcef0825886035c3091c8537a69535bc96f9e265f0ae45863543d40ae70090.
\end{center}
The finite language and its computation semantics are inherited unchanged
from package 19; its proof is included as \texttt{upstream/wang943\_language.md}.

Write $A_c=1024$.  In slot order west, north, east, south, the copy offsets
are $1023,1013,1003,983$.  The native support is
\begin{align*}
\mathcal E={}&[0,942]\cup[983,992]\cup[1003,1032]\\
 &{}\cup\{1063,1064,1095,1096,2195,3296\}.
\end{align*}
It has 989 elements.  If tile $i$ has edge labels $t_{i,0},\ldots,t_{i,3}$,
put $\epsilon_{i,s,j}$ equal to bit $j$ of $t_{i,s}+512$ and define
\[
 c_i=1+\sum_{s=0}^3\sum_{j=0}^9
          \epsilon_{i,s,j}A_c^{1+10s+j}+\sum_{j=41}^{44}A_c^j.
\]
The copy at exponent $\operatorname{offset}_s+j$ has coefficient
$A_c^{1+10s+j}$.  The anchors $1095,1096,2195,3296$ have coefficients
$A_c^{41},A_c^{42},A_c^{43},A_c^{44}$, respectively, and both dummy
coefficients are zero.  These rules specify every $c_e$.
The clause mask is
\[
 \mu=A_c-2+\sum_{j=1}^{44}(A_c-3)A_c^j+2^{450}(2^{78}-1).
\]
Set, now without five-power rounding,
\[
 b_c=529,\qquad L_c=12031,\qquad d_c=b_cL_c=6364399,
 \qquad \rho=2^{529},\qquad B=2^{6364399}.
\]
All the following quantities are fixed integers, not unknowns:
\begin{align*}
D_C={}&\rho^{30}+\rho^{5517}+\rho^{1100}+\rho^{2200}
       +\sum_{e\in\mathcal E}c_e
                 (\rho^{7631-e}+\rho^{8733-e}),\\
D_R={}&\rho^{5497},\\
M_{F,0}={}&\mu(\rho^{7631}+\rho^{8733})
       +\sum_{j=0}^9\rho^{1013+j}+\sum_{j=0}^9\rho^{6520+j}
       +\rho^{2195}+\rho^{3296},\\
M_C={}&B-1-\sum_{\substack{e\in\mathcal E\\e\ne1}}\rho^e
              -2\rho^{1063},\qquad M_{F,\mathrm{native}}=M_{F,0}+4.
\end{align*}
The four ports used by \emph{every} circuit in this delivery are
\begin{equation}\label{eq:four-ports}
 \boxed{\ \beta=B-1,\quad K_0=D_C+BD_R,\quad
 M_C,\quad M_F=M_{F,\mathrm{native}}+B-1.\ }
\end{equation}
The last port is the shifted high mask.  Replacing it by the unshifted
mask, or replacing only part of this tuple, defines a different and
unjustified instance.

\subsection{Why the smaller layout is sufficient}

The native clause coefficients are unchanged.  With
\[
 \mathcal M=(|\mathcal E|+2)\left(2\sum c_e+6\right),\qquad
 T_c=\max(4\mathcal M+8,2\mu+4,16),
\]
the exact check gives $\operatorname{bitlength}(T_c)=529$ and
$2^{529}\ge T_c$.  Thus the same uniform carry bound holds.
The largest unshifted product exponent is
$8733+3296=12029<12031$.  The earlier exact support proof of target
isolation, clean clause bands, and unique anchor alignment uses no
five-power property.  It remains valid for the new cell length.

In particular, the independent finite checks give
\begin{align*}
 \{(7631-e)\bmod L_c:e\in\mathcal E\}
 &\cap\{(8733-e)\bmod L_c:e\in\mathcal E\}=\varnothing,\\
 \{(2195-e)\bmod L_c:e\in\mathcal E\}
 &\cap\{(3296-e)\bmod L_c:e\in\mathcal E\}=\{0\}.
\end{align*}
All 69,230 native-basis convolution checks agree with the prescribed
overlap and clause contributions.  The population and residue conditions
are exactly
\begin{gather*}
0<M_C,M_{F,\mathrm{native}}<B-1,\qquad
M_C\equiv2\pmod4,\quad M_{F,\mathrm{native}}\equiv4\pmod8,\\
\operatorname{popcount}(M_C)+
\operatorname{popcount}(M_{F,\mathrm{native}})=d_c.
\end{gather*}
These calculations verify the finite hypotheses of the inherited
soundness proof.  That proof applies to arbitrary solutions and does not
assume that a recovered tableau size is a prime power.

\subsection{A fully explicit completeness controller}

Completeness needs a new argument: the old claim that the control
multiplier is a unit modulo a power of five is no longer applicable.
Factor
\[
 d_c=23^2\cdot53\cdot227,
 \qquad \lambda=452453,\qquad k_c=452452=\lambda-1.
\]
Trial division verifies that $\lambda$ is prime.  Exact modular
exponentiation gives
\begin{center}
\begin{tabular}{lrrrr}
\toprule
Modulus & $529$ & $53$ & $227$ & $452453$\\
\midrule
Order of $B$ & $11$ & $52$ & $226$ & $452452$\\
\bottomrule
\end{tabular}
\end{center}
Consequently $\operatorname{ord}_{d_c}(B)=64636$ divides $\lambda-1$,
and $B^{k_c}\equiv1\pmod{d_c}$.  Furthermore
\[
 B^{k_c}\bmod\lambda^2=155849698116\ne1,
 \qquad v_\lambda(B^{k_c}-1)=1.
\]
Put $H_0=1+BD_C+B^2(D_R+1)$.  The complete unit check is small:
\begin{center}
\begin{tabular}{rrrr}
\toprule
Prime $p$ & $B\bmod p$ & $(B^2-1)\bmod p$ & $H_0\bmod p$\\
\midrule
23 & 3 & 8 & 8\\
53 & 14 & 36 & 7\\
227 & 20 & 172 & 14\\
452453 & 167241 & 264979 & 406475\\
\bottomrule
\end{tabular}
\end{center}

Choose the spatial width $w_c$ and temporal height $h_c$ to be sufficiently
large powers of $\lambda$; independent rectangular padding permits this.
Write $N_c=w_ch_c=\lambda^n$, $q=B^{N_c}$, and impose
$N_c>d_c\lambda^2$ and $h_c>\lambda$.  Since all relevant orders divide
$\lambda-1$,
\[
 q\equiv B^{w_c}\equiv B\pmod{d_c\lambda}.
\]
The exact change in the packed Pell index when a lower dummy bit at
exponent $e_*=1063$ in cell $j$ is switched on is inherited unchanged:
\[
 \delta R=-\Gamma B^j,\qquad
 \Gamma=\rho^{e_*}(q^2-1)
        \bigl[1+q(D_C+BD_R+B^{w_c})\bigr].
\]
The table proves that $\Gamma$ is a unit modulo $d_cN_c$.

\begin{lemma}[Boolean coverage of the actual modulus]\label{lem:padding}
Every residue modulo $d_cN_c$ is the sum of a Boolean subset of
\[
 \{B^{2+k_cj}:0\le j<T\}\cup\{B^3\},\qquad T=N_c/\lambda.
\]
Every selected cell and its spatial and temporal translations used by the
compiler are nonwrapping.
\end{lemma}
\begin{proof}
Let $C_0=d_c\lambda$, so $d_cN_c=C_0T$ and $T>C_0$.
The valuation just proved and the binomial lifting identity show that
$B^{k_c}$ has order $T$ modulo $\lambda^n$, and order one modulo $d_c$.
Its first $T$ powers are therefore exactly the $T$ distinct residues
$1+C_0z$ modulo $C_0T$.

Normalize the desired residue by $B^{-2}$.  Choose
$\varepsilon\in\{0,1\}$ so that the normalized residue minus
$\varepsilon B$ is nonzero modulo $\lambda$.  Let $a$ be its least
residue modulo $C_0$.  Then $1\le a<C_0<T$, $\gcd(a,T)=1$, and the
remaining target has the form $a+C_0u$.  Choose $z_0$ solving
\[
 az_0+\frac{a(a-1)}2\equiv u\pmod T.
\]
Take the $a$ distinct orbit entries with additive coordinates
$z_0,z_0+1,\ldots,z_0+a-1$ modulo $T$, and include $B^3$ exactly when
$\varepsilon=1$.  Their sum is the target.  No position is reused.

The largest available position is
$2+(\lambda-1)(N_c/\lambda-1)=N_c-N_c/\lambda-\lambda+3$.
Since $h_c>\lambda$ implies $w_c<N_c/\lambda$, its translates by one
and by $w_c$ stay below $N_c$.  The extra position 3 is distinct because
$k_c>1$ and satisfies the same conditions.
\end{proof}

For a canonical blank halting tableau, begin with all these dummy bits
zero and let its actual packed index be $R_0$.  Apply the lemma to
\[
 (R_0-d_cw_c)\Gamma^{-1}\pmod{d_cN_c}.
\]
The modified tableau has
$R\equiv d_cw_c\pmod{d_cN_c}$, exactly the temporal rotation required
by the power $X=2^R$.  The selected dummy bits change no tile, overlap,
marker, or carry condition.  The unchanged carry estimate preserves
the strict slack $2C+F<q/3$ and hence positive transport coordinates.
The inherited Pell converse supplies the remaining witnesses.

All dimensions $b_c,L_c,d_c,w_c,h_c,N_c$ are odd.  More explicitly, write
$u_c=d_cN_c=\log_2q$.  The unchanged low-mask congruences give
$R\equiv3\pmod4$, so the canonical choice $w=2^{R-u_c}$ is a positive
integer square.  For the canonical even Pell parameter $A_c'=a+2$,
the coordinate $c=\psi_{A_c'}(R)$ is odd.  Choosing the strong rank
$m_c=cR$ makes $i=\psi_{A_c'}(m_c)/c^2$ an odd positive integer, by
the binomial expansion in the signed source \cite{normalized,oddscale}.
The exact population identity still gives
$\operatorname{popcount}(R)=3u_c+2$, and the half-binomial identity gives
$v_2(Y)=3u_c$; hence the scale $s=Y/q^3$ is odd.
This establishes the full adapter needed for the 24-integer construction;
the modular checks alone would not establish it.

\begin{theorem}[Complete coefficient replacement]\label{thm:coefficients}
The complete tuple \eqref{eq:four-ports} satisfies the same blank-K215
native compiler interface as package 19, including the signed source's
canonical odd and square witness requirements.  It may replace that
tuple in the compatible arithmetic forms in this report.
\end{theorem}

\subsection{Actual bytes and the size improvement}

The four integers are included as gzip-compressed, minimal unsigned
little-endian binary payloads.  They are already materialized, not merely
specified by a program waiting for K215 to halt.
\begin{center}
\begin{tabular}{lrrr}
\toprule
Port & Bits & Raw bytes & Compressed bytes\\
\midrule
$\beta$ & 6,364,399 & 795,550 & 806\\
$K_0$ & 9,272,313 & 1,159,040 & 15,461\\
$M_C$ & 6,364,399 & 795,550 & 1,151\\
$M_F$ & 6,364,400 & 795,550 & 901\\
\midrule
Total & & 3,545,690 & 18,319\\
\bottomrule
\end{tabular}
\end{center}
Package 19 used $d_c=9,765,625$ and 5,312,269 raw bytes.  The new $d_c$
is smaller by about $34.83\%$, and the raw payload by about $33.25\%$.
All four fixed ports have smaller numerical magnitudes.  Compression size
is reported separately; it is not an algebraic coefficient-height measure.
The values 529 and 12031 are the least odd choices satisfying the existing
carry threshold and existing maximum support bound, respectively.  This
limited minimality claim fixes the tile language, mask and layout; it is
not a minimum over all coefficient encodings.

% END sections/coefficients.tex

% BEGIN sections/reproduction.tex
\section{Reproducibility, limitations, and remaining directions}
\label{sec:reproduction}

\subsection{What is delivered}

The archive contains the standalone \LaTeX{} source of this report and its
compiled PDF, the editable section sources, twenty literal arithmetic
circuits, all four compressed coefficient integers, standard-library
Python builders and checkers, and the pinned upstream proof and machine
files used in the composition.  No repository checkout or network
connection is needed to run the numerical and algebraic verification.
Upstream files retain their source paths, commit identifiers, checksums,
and notices.  The inherited formal endpoint is a source record, not a
self-contained copy of the whole Lean project.

For a practical entry point, use
\texttt{signed\_separating\_saturated24\_product.json} for unrestricted
integers, or
\texttt{optimized\_tightest\_separating12\_product.json} for the
degree-1069 positive equation.  Replacing \texttt{product} with
\texttt{system} selects the corresponding three-equation version.  Every
circuit explicitly refers to the one delivered coefficient binding.  The
binding records minimal unsigned little-endian bytes and SHA-256 hashes;
\texttt{MF} always means the shifted high mask.

From the unpacked archive directory, the standard-library verification
and reconstruction entry points are
\begin{verbatim}
python3 verify.py
python3 reproduce.py
pdflatex -interaction=nonstopmode -halt-on-error con_zfc_synthesis.tex
pdflatex -interaction=nonstopmode -halt-on-error con_zfc_synthesis.tex
\end{verbatim}
The first command checks the supplied artifacts.  The second regenerates
the new coefficient tuple and the four principal construction families,
then rebinds and verifies the delivered circuits.  The six exact-refactored
comparison circuits and the two inherited comparison circuits are included
literally and checked against their preserved source packets; they do not
depend on a missing external generator.  The final two commands rebuild
the document.  The Python scripts require no third-party Python package.
The \LaTeX{} source uses ordinary TeX Live packages.

\subsection{The verification boundary}

The finite audits check the literal machine-to-tile enumeration, the
layout's support and carry inequalities, the modular padding constants,
the four materialized payloads, every circuit's well-formedness and live
inputs, exact small polynomial identities, and independent evaluations
of the norm constructions.  Full univariate finite-field calculations
use the actual four coefficient residues to check the principal degree
claims.  The positive tight-weight and signed-variable sections also
give uniform nonzero leading forms, so the claimed degrees are not based
only on random evaluations or generic syntactic upper bounds.

These checks have specific limits.  Finite-field agreement is a
diagnostic unless backed by the displayed polynomial identity or degree
proof.  A successful program run does not prove universal compiler
soundness, does not formally verify the root-recovery arguments, and does
not establish the truth of consistency.  The infinite claims rest on
the mathematical proofs and their stated inherited interfaces.  This
report makes no new end-to-end Lean or Rocq verification claim.

The supplied collection is itself a set of research reports; its intake
note explicitly distinguishes claims in the packages from independently
re-derived facts \cite{collection}.  This synthesis follows that
distinction: it preserves the exact sources for review and supplies new
proofs for the changed constructions, rather than treating comparison
metadata as a substitute for the reduction.

\subsection{An attempted extra variable reduction that is not claimed}

The remaining index quotient $h$ is tempting to merge with the main-root
quotient.  A single combined carrier, however, supplies a combined
divisibility condition.  Recovering separately the index congruence
modulo $E$ and the main-root congruence modulo $H$ would require further
arithmetic information, such as an appropriate coprimality condition.
There is no available identity forcing $\gcd(E,H)=1$:
\[
 E=XY,\qquad H=4(E+Y)+3,
\]
and the elementary parameter restrictions do not justify that inference.
The required decomposition cannot simply be assumed.  Consequently no
eleven-positive-variable theorem is asserted here.  A valid further
reduction would have to preserve both congruences, or prove an additional
property of every recovered Pell witness before using it.

There are still concrete directions for improvement.  The 943-tile
language and its spatial exponent layout were kept fixed, so reducing
the local alphabet or redesigning the layout could lower coefficient
height further.  A different norm elimination might improve the tradeoff
between twelve-variable degree and arithmetic length.  Fewer integer
coordinates would require a replacement for one of the six positivity
guards, another valid root merger, or a stronger signed interface.
None of these potential improvements is counted in the delivered results.


% END sections/reproduction.tex

% BEGIN sections/references.tex
\clearpage
\begin{thebibliography}{99}
\bibitem{collection}
ProveIt. \emph{Con(ZFC) in Diophantine form: wave-3 reconciliation}.
\href{https://github.com/VladimirReshetnikov/ProveIt/blob/ae3007ca4467011e0e051c84acb4a1dd13971b0e/SetTheory/Cardinals/docs/reports/hilbert-tenth-problem/con-zfc-diophantine-forms/README.md}{Pinned repository source} at \texttt{ae3007ca4467}. Local copy: \path{upstream/collection_at_pin.md}.

\bibitem{repo18}
ProveIt. \emph{Package 18: Three explicit Diophantine forms of ZFC consistency}.
\href{https://github.com/VladimirReshetnikov/ProveIt/blob/ae3007ca4467011e0e051c84acb4a1dd13971b0e/SetTheory/Cardinals/docs/reports/hilbert-tenth-problem/con-zfc-diophantine-forms/18-jones9/README.md}{Pinned repository source} at \texttt{ae3007ca4467}.

\bibitem{repo19}
ProveIt. \emph{Package 19: Twelve positive variables and materialized coefficients}.
\href{https://github.com/VladimirReshetnikov/ProveIt/blob/ae3007ca4467011e0e051c84acb4a1dd13971b0e/SetTheory/Cardinals/docs/reports/hilbert-tenth-problem/con-zfc-diophantine-forms/19-wang943/README.md}{Pinned repository source} at \texttt{ae3007ca4467}.

\bibitem{repo20}
ProveIt. \emph{Package 20: Twelve positive unknowns and materialized coefficients}.
\href{https://github.com/VladimirReshetnikov/ProveIt/blob/ae3007ca4467011e0e051c84acb4a1dd13971b0e/SetTheory/Cardinals/docs/reports/hilbert-tenth-problem/con-zfc-diophantine-forms/20-root12/README.md}{Pinned repository source} at \texttt{ae3007ca4467}.

\bibitem{repo21}
ProveIt. \emph{Package 21: A twelve-variable equation for Con(ZFC)}.
\href{https://github.com/VladimirReshetnikov/ProveIt/blob/ae3007ca4467011e0e051c84acb4a1dd13971b0e/SetTheory/Cardinals/docs/reports/hilbert-tenth-problem/con-zfc-diophantine-forms/21-norm12/proofs/NORM12_PROOF.md}{Pinned repository source} at \texttt{ae3007ca4467}. Local copy: \path{upstream/norm12_proof.md}.

\bibitem{repo22}
ProveIt. \emph{Package 22: Monic merger proof}.
\href{https://github.com/VladimirReshetnikov/ProveIt/blob/ae3007ca4467011e0e051c84acb4a1dd13971b0e/SetTheory/Cardinals/docs/reports/hilbert-tenth-problem/con-zfc-diophantine-forms/22-monic12/proofs/merger_proof.md}{Pinned repository source} at \texttt{ae3007ca4467}. Local copy: \path{upstream/monic_merger_proof.md}.

\bibitem{repo23}
ProveIt. \emph{Package 23: Separating-root twelve-variable construction}.
\href{https://github.com/VladimirReshetnikov/ProveIt/blob/ae3007ca4467011e0e051c84acb4a1dd13971b0e/SetTheory/Cardinals/docs/reports/hilbert-tenth-problem/con-zfc-diophantine-forms/23-separating12/README.md}{Pinned repository source} at \texttt{ae3007ca4467}.

\bibitem{K215}
ProveIt. \emph{KEndpoint215: blank halting and ZFC refutation}.
\href{https://github.com/VladimirReshetnikov/ProveIt/blob/ae3007ca4467011e0e051c84acb4a1dd13971b0e/Computability/BusyBeaver/Lean/BusyBeaver/ZFC/KEndpoint215.lean}{Pinned repository source} at \texttt{ae3007ca4467}. Local copy: \path{upstream/KEndpoint215.lean}.

\bibitem{baseline67}
ProveIt. \emph{Fixed blank 67-operation polynomial and native compiler interface}.
\href{https://github.com/VladimirReshetnikov/ProveIt/blob/ae3007ca4467011e0e051c84acb4a1dd13971b0e/SetTheory/Cardinals/docs/reports/hilbert-tenth-problem/con-zfc-diophantine-forms/21-norm12/upstream/FIXED_BLANK_67_PROOF.md}{Pinned repository source} at \texttt{ae3007ca4467}. Local copy: \path{upstream/fixed_blank_67_proof.md}.

\bibitem{normalized}
ProveIt. \emph{Signed-domain proof for the normalized Pell interface}.
\href{https://github.com/VladimirReshetnikov/ProveIt/blob/ae3007ca4467011e0e051c84acb4a1dd13971b0e/SetTheory/Cardinals/docs/reports/hilbert-tenth-problem/con-zfc-diophantine-forms/21-norm12/upstream/SIGNED_DOMAIN_29.md}{Pinned repository source} at \texttt{ae3007ca4467}. Local copy: \path{upstream/signed_domain29.md}. The ordinary and helical compiler proofs are also included as \path{upstream/ordinary_blank_proof.md} and \path{upstream/blank_helical_proof.md}.

\bibitem{oddscale}
ProveIt. \emph{Odd-scale 27-integer proof and negative-scale exclusion}.
\href{https://github.com/VladimirReshetnikov/ProveIt/blob/ae3007ca4467011e0e051c84acb4a1dd13971b0e/SetTheory/Cardinals/docs/reports/hilbert-tenth-problem/con-zfc-diophantine-forms/23-separating12/proofs/ODD_SCALE_27_PROOF.md}{Pinned repository source} at \texttt{ae3007ca4467}. Local copy: \path{upstream/odd_scale27.md}.

\bibitem{wang}
ProveIt. \emph{Wang943 coefficient language and compact-layout proof}.
\href{https://github.com/VladimirReshetnikov/ProveIt/blob/ae3007ca4467011e0e051c84acb4a1dd13971b0e/SetTheory/Cardinals/docs/reports/hilbert-tenth-problem/con-zfc-diophantine-forms/19-wang943/compiler/WANG943_COMPACT_LAYOUT.md}{Pinned repository source} at \texttt{ae3007ca4467}. Local copy: \path{upstream/wang943_compact_layout.md}. The separate language proof is included as \path{upstream/wang943_language.md}.

\bibitem{padding}
ProveIt. \emph{Pell-kernel prime padding with an arbitrary fixed odd factor}.
\href{https://github.com/VladimirReshetnikov/ProveIt/blob/ae3007ca4467011e0e051c84acb4a1dd13971b0e/Computability/HilbertTenthProblem/Papers/research-wip/native-stream-queue/pell_kernel_prime_padding.md}{Pinned repository source} at \texttt{ae3007ca4467}. Local copy: \path{upstream/prime_padding.md}.

\bibitem{retained59}
ProveIt. \emph{Retained-variable 59-operation system}.
\href{https://github.com/VladimirReshetnikov/ProveIt/blob/ae3007ca4467011e0e051c84acb4a1dd13971b0e/SetTheory/Cardinals/docs/reports/hilbert-tenth-problem/con-zfc-diophantine-forms/22-monic12/proofs/RETAINED59_PROOF.md}{Pinned repository source} at \texttt{ae3007ca4467}. Local copy: \path{upstream/retained59_proof.md}.

\end{thebibliography}

% END sections/references.tex

\end{document}
